\clearpage\begingroup
\HistoricalHead{149}{R}{GAUSS, SCHUMACHER, OLBERS.}
\begin{center}
[7.]\par\medskip
\textbf{\textsc{Olbers} an \textsc{Gauss}.\quad Bremen, 10. Juli 1813.}\par\medskip
\rule{.80\linewidth}{.4pt}\par
[\textsc{Wilhelm Olbers}, Sein Leben und seine Werke II\,1, S. 524.]\par
\rule{.80\linewidth}{.4pt}
\end{center}
\noindent\{. . . . . . . . . . . . . . . . . . . . .\par
\noindent\hspace*{\parindent}Bei dieser Gelegenheit bitte ich Sie um eine Belehrung. Es scheint mir, dass wir die Aperturen unserer nicht achromatischen Cometensucher viel zu gross machen, wodurch die Bestimmtheit und Schärfe des Bildes leidet, ohne dass wir an Licht gewinnen. Dass die Lichtstärke eines Fernrohrs bei gleicher Durchsichtigkeit der Gläser wie das Quadrat des Durchmessers des Objectiv-Glases sich verhält, ist nur so lange wahr, als der Strahlenbüschel, der aus
\clearpage\SourceLeaf{155}\HistoricalHead{150}{L}{PHYSIK.\quad BRIEFWECHSEL.}
\noindent dem letzten Augenglase kömmt, im Durchmesser kleiner, als der Durchmesser der Pupille unseres Auges ist. Dieser Durchmesser unserer Pupille wird auch des Nachts nicht viel über $2\frac12$ Pariser Linien betragen. Vergrössert nun ein Cometensucher $m$mal, und ist der Durchmesser unserer Pupille $=d$, so ist $md$ die einzig nützliche Apertur des Cometensuchers. Also für 8malige Vergrösserung und $d=2\frac12$ Linien gesetzt, würde $md=20$ Linien seyn. Ist die Apertur grösser, so gewinnt man nichts an Licht, nur der Durchmesser des Strahlenbüschels wird grösser, als ihn die Apertur unserer Pupille fassen kann. Es scheint mir also, dass wir die Objective unserer Cometensucher nie breiter als 2 Zoll, oder wenn wir 10mal vergrössern wollen, höchstens $2\frac12$ Zoll machen sollten. Die von \textsc{Ramsden} sind viel breiter. . . . . . .\}
\bigskip\begin{center}
[8.]\par\medskip
\textbf{\textsc{Gauss} an \textsc{Olbers}.\quad Göttingen, 13. September 1813.}\par\medskip
\rule{.80\linewidth}{.4pt}\par
[\textsc{Wilhelm Olbers}, Sein Leben und seine Werke II\,1, S. 530.]\par
\rule{.80\linewidth}{.4pt}
\end{center}
\noindent . . . . . . . . . . . . . . . . . . . . .\par
\noindent\hspace*{\parindent}Ihre Bemerkung wegen der Objective schwach vergrössernder Fernröhre ist sehr gegründet. Von zwei übrigens gleichen Cometensuchern, wo bei einem das Objectiv nur eben äqual dem Product ist aus der Vergrösserung in die grösste Öffnung der Pupille, bei dem andern aber grösser, hat das erstere den entschiedenen Vorzug grösserer \textls[120]{Leichtigkeit} und Wohlfeilheit, ohne dass das andere mehr leistete, wenn anders \textls[120]{das Auge genau an seinem Platze steht.} Sobald aber diese Bedingung aufhört, hat man bei letzterem zwar etwas weniger Schärfe, aber keinen Verlust an Helligkeit, was bei ersterem sich umgekehrt verhält. Ich glaube daher, dass es doch wol gut ist, die Öffnung des Objectivs \textls[120]{etwas} grösser zu machen, als jenes Product, aber doch nicht so gross, wie bei den meisten Kometensuchern, die ich unter Händen gehabt habe. . . . . . .
\clearpage\SourceLeaf{156}\HistoricalHead{151}{R}{GAUSS AN SCHUMACHER UND ENCKE.}
\begin{center}
[9.]\par\medskip
\textbf{\textsc{Gauss} an \textsc{Schumacher}.\quad Göttingen, 7. Juni 1836.}\par\medskip
\rule{.75\linewidth}{.4pt}\par
[Briefwechsel zwischen \textsc{Gauss} und \textsc{Schumacher} III, S. 62.]\par
\rule{.75\linewidth}{.4pt}
\end{center}
\noindent . . . . . . . . . . . . . . . . . . . . .\par
\noindent\hspace*{\parindent}Ich habe mich in der letzten Woche viel mit optischen, namentlich mit den sehr interessanten Beugungsversuchen beschäftigt. \textsc{Schwerdt} hat in diesem Felde recht viel geleistet[*)], um diese höchst mannigfaltigen Erscheinungen, wozu \textsc{Fraunhofer} zuerst den Weg öffnete, aus Einem Princip abzuleiten. Inzwischen bleibt doch noch sehr viel übrig, bis die Theorie als vollständig und erschöpfend angesehen werden kann.
\begingroup\renewcommand{\thefootnote}{}\footnotetext{[*) F. M. \textsc{Schwerd}, \textit{Die Beugungserscheinungen aus den Fundamentalgesetzen der Undulationstheorie analytisch entwickelt u.s.w.}, Mannheim 1835.]}\endgroup
\bigskip\begin{center}
[10.]\par\medskip
\textbf{\textsc{Gauss} an \textsc{Encke}.\quad Göttingen, 3. Juni 1836.}
\end{center}
\noindent . . . . . . . . . . . . . . . . . . . . .\par
\noindent\hspace*{\parindent}In den letzten Wochen habe ich angefangen, mich mit einigen optischen Untersuchungen, namentlich auch mit den Diffractionsversuchen, etwas zu beschäftigen, welche reiche Gegenstände mich fast ebenso sehr zu interessiren anfangen, wie die magnetischen und elektromagnetischen. Beschäftigung mit den letzteren hat inzwischen dabei ihren Fortgang und wir werden in kurzem die interessanten Versuche über die Induction durch den blossen Erdmagnetismus, welche wir schon im März 1835 mit schwächern und im März 1836 mit stärkeren [Apparaten] angestellt hatten, mit noch viel stärkeren zu wiederholen [haben,] um alles auf genaue Zahlenbestimmungen zurückbringen zu können. Bei diesen Versuchen und den damit zusammenhängenden wird, in meinem Inductor, meinem neuen Multiplicator und in einem grossen Inductionsrade zusammen, \textls[120]{eine Meile} übersponnenen Drahts gebraucht werden.\par
\noindent . . . . . . . . . . . . . . . . . . . . .
\clearpage\endgroup
